% !TEX root = ../../hw01.tex
% Homework 01, Exercise 10. Shared solution.

\begin{proof}[Solution]
Since each $A_n$ is countably infinite, there exists a bijection
from $\N$ to $A_n$. Using the axiom of countable choice, select
one such bijection for each $n$. This gives a list without
repetition within each row, which we write as
\[
  A_n=\{a_{n,1},a_{n,2},a_{n,3},\ldots\}.
\]
Put these lists into rows. Reading one whole row first would
never reach the next row. Instead, follow successive diagonals
$n+m=2,3,\ldots$, alternating direction as in
Figure~\ref{fig:hw01:diagonal-enumeration}. This gives
\[
  a_{1,1},\ a_{1,2},\ a_{2,1},\ a_{3,1},
  \ a_{2,2},\ a_{1,3},\ a_{1,4},\ldots.
\]

\begin{figure}[h]
  \centering
  % !TEX root = ../../problems/hw01.tex
\begin{tikzpicture}[x=1.85cm,y=0.72cm,>=Stealth,font=\small]
  \foreach \col in {1,2,3,4} {
    \node[font=\footnotesize] at (\col,0.75) {$m=\col$};
    \node at (\col,-4) {$\vdots$};
  }
  \foreach \row in {1,2,3,4} {
    \node[anchor=east] at (0.4,{1-\row}) {$A_{\row}$};
    \node at (4.75,{1-\row}) {$\cdots$};
    \foreach \col in {1,2,3,4}
      \node (a\row\col) at (\col,{1-\row}) {$a_{\row,\col}$};
  }
  \node at (4.75,-4) {$\ddots$};
  \foreach \from/\to in {11/12,12/21,21/31,31/22,22/13,
    13/14,14/23,23/32,32/41} {
    \draw[->,black!60,line width=0.85pt,shorten <=2pt,shorten >=2pt]
      (a\from) -- (a\to);
  }
  \draw[->,black!60,line width=0.85pt]
    (a41.south) -- (1,-3.7);
\end{tikzpicture}

  \caption[Enumerating a countable union.]{Follow the arrows along diagonals of constant $n+m$.
    Every entry is reached after finitely many steps.}
  \label{fig:hw01:diagonal-enumeration}
\end{figure}

Each diagonal is finite, and each $a_{n,m}$ has only finitely
many diagonals before its own. Hence every entry is eventually
reached. Every element of $A$ occurs in some row, so this sequence
contains all of $A$.

The sets may overlap. Skip any value already listed, keeping only
its first occurrence. Every element has a first occurrence at
a finite position, so none is lost by this procedure. The
retained list contains each element of $A$ exactly once.
Since $A_1\subseteq A$ is infinite, the retained list is also
infinite. Numbering its entries in order gives a bijection
from $\N$ to $A$. Therefore
\[
  |A|=\aleph_0.\qedhere
\]
\end{proof}
